\paragraph{The reverse mass-twelve row with a mass-four action.}
We now exclude the remaining orientation
\[
 (w,g_1,g_2,d_1,d_2,u_1,u_2)
                  =(12,96,80,4,12,12,4).
\]
Write $G_+=G_1$, $G_-=G_2$, $U_+=U_1$, $U_-=U_2=D$,
$\delta_+=\psi=D(D-3)$, and $\delta_-=\psi+2R$.
Here $R$ is an effective symmetric character of mass four.
In particular $|\delta_-(0)|\le8$ and $|U_+(0)|\le8$.

Suppose first $W(0)\ne0$, so $W(0)=\pm8$.
The minus-cell identity
$DG_+=W(G_--D\delta_-)$ and coprimality of $W,D$
give $G_+=WQ$ with $Q$ integral, including by the content
argument if $W$ is nonprimitive.
Then $G_-=D(Q+\delta_-)$.
The plus-cell identity at zero implies $Q(0)=0$, because
its left side is $U_+(0)G_-(0)=0$ and its right side is
$W(0)^2Q(0)$.
Thus $Q=DC$ and
\[
 G_+=DWC,\qquad G_-=D\delta_-+D^2C,\qquad C(4)=2.
\]
The trace parity $G_1\equiv0$, together with $W\equiv D$
modulo two, forces $C\equiv0\pmod2$.
Dividing the plus-cell identity by $D$ now gives
\[
 U_+(\delta_-+DC)=W(WC-\psi).
\]
At zero this says
$64C(0)=U_+(0)\delta_-(0)$.
The right side has absolute value at most sixty-four.
Therefore $|C(0)|\le1$; its even parity forces $C(0)=0$.
The integral polynomial $C$ would be divisible by $D$,
contradicting $C(4)=2$.

Suppose next $W(0)=0$, so $W=DH$ with the previously
proved mass-twelve classification. Since $H(4)=3$ and
$G_-(4)=80$, the same irreducibility/content argument
as in the preceding mass-twelve rows forces $U_+=W$.
More explicitly, a nonconstant $H$ is primitive irreducible
and cannot divide $G_-$; hence it divides $U_+$.
Its odd value $H(0)\ne\pm1$ forces $U_+(0)=0$,
and the classification of $U_+/D$ gives $U_+=W$.
When $H=3$, reduction modulo three and the mass-twelve
effectivity/parity bound give the same conclusion.
Consequently
\[
 G_+=G_-+D\psi,\qquad
                  (H-1)G_-=D(H\delta_-+\psi).
\]
The necessary NN-G030 congruence is
$G_--D\delta_-\equiv G_+\pmod4$.
It forces $R\equiv\psi\pmod2$.
The effective character $R$ has mass four, while
$\psi=T_1+T_2$ already has four odd Laurent coefficients
and mass four. Hence coefficientwise effectivity and
the mass bound force $R=\psi$ exactly. In particular
$\delta_-=3\psi$ and
\[
                    (H-1)G_-=D\psi(3H+1).
\]
Put $J=(H-1)/2\in\mathbb Z[D]$.
Then $J(4)=1$. It is either the constant one (when $H=3$)
or a monic nonconstant polynomial, since the leading
coefficient of a nonconstant $H$ is two.
The preceding identity becomes
\[
                    J(G_--3D\psi)=2D\psi.
\]
Thus the monic $J$ divides $2D^2(D-3)$ in $\mathbb Z[D]$.
Its only possibilities over $\mathbb Q[D]$ are
$D^a(D-3)^b$ with $0\le a\le2$, $0\le b\le1$;
monicity fixes the scalar. Its value one at four
excludes $a>0$. Therefore $J=1$ or $J=D-3$.

If $J=1$, the identities force
\[
 W=3D,\quad G_-=5D\psi,\quad G_+=6D\psi,\quad
 S=10D^3(D-3)^2.
\]
If $J=D-3$, they force
\[
 W=D(2D-5),\quad G_-=D^2(3D-7),\quad
 G_+=2D^2(2D-5),\quad S=2D^3(3D-7).
\]
In each case $S=G_+G_-/W$ is the actual bottom
structure coefficient. Its prescribed coefficientwise
lower bound is $4D^2(D-1)^2$, whose central Laurent
coefficient is $144$.
Using $[q^0]D^h=\binom{2h}{h}$, the two forced
central coefficients of $S$ are respectively
\[
 10(252-6\cdot70+9\cdot20)=120,\qquad
                    6\cdot70-14\cdot20=140.
\]
Both are smaller than $144$, contradicting effectivity
and the given trace domination. This excludes the reverse
row and every permitted auxiliary increment, without a
support-degree restriction.
\hfill$\square$

This is an integral-character obstruction in the same
closed kernel presentation as NN-G036. The final inequality
compares actual multiplicities of the trivial torus character;
the dictionary is SGA~3, I, 4.7.2--4.7.3 and NN-G010.
No conclusion about a different multiplication table follows.
